% Part: normal-modal-logic
% Chapter: syntax-and-semantics
% Section: relational-models

\documentclass[../../../include/open-logic-section]{subfiles}

\begin{document}

\olfileid[pt]{nml}{syn}{rel}

\olsection{Modelos relacionais}

O conceito fundamental da semântica para lógicas modais normais é o de
\emph{modelo relacional}. Ele consiste em um conjunto de mundos,
relacionados por uma ``relação de acessibilidade'' binária, juntamente
com uma atribuição que determina quais !!{propositional variable}s
contam como ``verdadeiras'' em quais mundos.

\begin{defn}
  Um \emph{modelo} para a linguagem modal básica é uma tripla $\mModel{M}
  = \tuple{W, R, V}$, em que
  \begin{enumerate}
  \item $W$ é um conjunto não vazio de ``mundos'';
  \item $R$ é uma relação binária de acessibilidade sobre~$W$; e
  \item $V$ é uma função que atribui a cada !!{propositional
    variable}~$p$ um conjunto $V(p)$ de mundos possíveis.
  \end{enumerate}
  Quando vale $Rww'$, dizemos que $w'$ é \emph{acessível a partir
    de}~$w$. Quando $w \in V(p)$, dizemos que $p$ é \emph{verdadeira
    em}~$w$.
\end{defn}

A grande vantagem da semântica relacional é que os modelos podem ser
representados por diagramas simples, como o de \olref{fig:simple}. Os
mundos são representados por nós, e o mundo~$w'$ é acessível a partir
de~$w$ precisamente quando há uma seta de~$w$ para~$w'$. Além disso,
rotulamos um nó (mundo) com~$\mTrue{p}$ quando $w \in V(p)$ e, caso
contrário, com~\mFalse{p}. \olref{fig:simple} representa o modelo com
$W = \{w_1, w_2, w_3\}$, $R = \{\tuple{w_1, w_2},\tuple{w_1, w_3}\}$,
$V(p) = \{w_1, w_2\}$ e $V(q) = \{w_2\}$.

\begin{figure}
  \begin{center}
    \begin{tikzpicture}[modal]
      \node[world] (w1) [label={[align=right]right:\mTrue{p}\\ \mFalse{q}}]{$w_1$}; 
      \node[world] (w2) [label={[align=right]right:\mTrue{p}\\ \mTrue{q}},
        above right=of w1]{$w_2$}; 
      \node[world] (w3) [label={[align=right]right:\mFalse{p}\\ \mFalse{q}},
        below right=of w1] {$w_3$};
      \draw[->] (w1) to (w2);
      \draw[->] (w1) to (w3);
    \end{tikzpicture}
  \end{center}
  \caption{Um modelo simples.}
  \ollabel{fig:simple}
\end{figure}

\end{document}
